\documentclass{ximera}
\title{Exponential Functions CW}

\newcommand{\pskip}{\vskip 0.1 in}

\begin{document}
\begin{abstract}
Exponential problems.
\end{abstract}
\maketitle


\section{Conceptual Questions}

\begin{question}  \label{Q39fvc3m3}
At some moment the balance in an account is increasing at the relative rate of $4\%$/yr. 
Use the language of \emph{specific} small changes to interpet the meaning of this statement. 
\end{question}


\begin{question}  \label{Qomfe03r3}

Suppose one population is always twice another. What can you say about their
\begin{enumerate}

\item average growth rates over any time interval?

\item growth rates at any moment in time?

\item relative changes over any interval of time?

\item relative average growth rates over any interval of time?

\item relative growth rates at any moment in time?

\end{enumerate} 

\item Answer (a)-(e) if instead one population were always 10,000 greater than the other.

\end{question}


\begin{question} \label{Qif9f3r3r3fr}
The function $G = f(v)$ \, , \, $20 \leq v \leq90$, expresses the gas mileage of a car (measured in miles/gal) in terms of its speed (measured in miles/hr).

Suppose
\[
 \left( \frac{1}{G} \cdot \frac{dG}{dv}  \right) \Big|_{v=70} = -0.05.
\]

\begin{enumerate}
\item What are the units of the expression above?

\item Explain its meaning in terms of \emph{specific} small changes.
\end{enumerate}
\end{question}

\begin{question} \label{Qpdov34r34rr3}
Between trip odometer readings $s=90$ and $s=100$ miles, the dashboard reading indicates you always have $50$ miles left to drive before running out of gas.

\begin{enumerate}
\item Find the instantaneous relative rate of change in the volume of gas in your tank during this period of the trip.

\item (computation) Find the relative change in the volume of gas in your tank during the 10-mile stretch from odometer reading $s=90$ miles to odometer reading $s=100$ miles.

\end{enumerate}
\end{question}


\section{Computational Questions}

\begin{question} \label{Qpvem3r34}
The function
\[
        P = 250 (1.1)^t \, , \, 0\leq t \leq 10,
\]
expresses the population of a colony of bacteria (in millions) in terms of the number of hours past noon.

\begin{enumerate}
\item Show algebraically that the relative average rate of change in the population from time $t=a$ to time $t=b>a$ hours past noon depends only on the difference $b-a$.

\item Use your expression from part (a) to write the relative instantaneous growth rate of the population as a limit.

\item Interpret your expresssion from part (b) geometrically as the slope of some tangent line.

\item Use numerical methods to approximate the relative instantaneous growth rate.

\item Find the exact relative instantaneous growth rate.

\end{enumerate}
\end{question}


\begin{question} \label{Qeirefr}
The function 
\[
     P = f(t) = 50 e^{\frac{t}{10}} \, , \, 0\leq t \leq 20 ,
\]
expresses the population (in millions) of a colony of bacteria in terms of the number of hours past  midnight. Graph shown below.

\begin{enumerate}
\item Find a function $r=g(P)$ that expresses the growth rate (measured in millions of bacteria/hour) in terms of the popuation (measured in millions). Include a domain.

\item Find the \emph{exact} growth rate when there are $200$ million bacteria. No calculator.

\item What is the population when it is increasing at the rate of $8$ million bacteria/hour?
\end{enumerate}
\end{question}



\end{document}